\documentclass[12pt,a4paper,oneside]{report}


\usepackage{amsmath,amssymb,calc,ifthen}

\usepackage{float}
%\usepackage{cancel}

\usepackage[table,usenames,dvipsnames]{xcolor} % for coloured cells in tables

\usepackage{tikz}

\usepackage{mathptmx}

\usepackage{hyperref}
\hypersetup{
    colorlinks,
    citecolor=black,
    filecolor=black,
    linkcolor=black,
    urlcolor=black
}

% Nice package for plotting graphs
% See excellent guide:
% http://www.tug.org/TUGboat/tb31-1/tb97wright-pgfplots.pdf
\usetikzlibrary{plotmarks}
\usepackage{amsmath,graphicx}
\usepackage{epstopdf}
\usepackage{caption}
\usepackage{subcaption}

% highlight - useful for TODOs and similar
\usepackage{color}
\newcommand{\hilight}[1]{\colorbox{yellow}{#1}}

\newcommand\ci{\perp\!\!\!\perp} % perpendicular sign
\newcommand*\rfrac[2]{{}^{#1}\!/_{#2}} % diagonal fraction

\usepackage{listings}



% margin size
\usepackage[margin=2.5cm]{geometry}

\tikzstyle{state}=[circle,thick,draw=black, align=center, minimum size=2.1cm,
inner sep=0]
\tikzstyle{vertex}=[circle,thick,draw=black]
\tikzstyle{terminal}=[rectangle,thick,draw=black]
\tikzstyle{edge} = [draw,thick]
\tikzstyle{lo} = [edge,dotted]
\tikzstyle{hi} = [edge]
\tikzstyle{trans} = [edge,->]


\definecolor{mygreen}{rgb}{0,0.6,0}
\definecolor{mygray}{rgb}{0.5,0.5,0.5}
\definecolor{mymauve}{rgb}{0.58,0,0.82}

\lstset{ %
  backgroundcolor=\color{white},   % choose the background color; you must add 
%\usepackage{color} or \usepackage{xcolor}
  basicstyle=\footnotesize,        % the size of the fonts that are used for the 
%code
  breakatwhitespace=false,         % sets if automatic breaks should only happen 
%at whitespace
  breaklines=true,                 % sets automatic line breaking
  captionpos=b,                    % sets the caption-position to bottom
  commentstyle=\color{mygreen},    % comment style
  deletekeywords={...},            % if you want to delete keywords from the 
%given language
  escapeinside={\%*}{*)},          % if you want to add LaTeX within your code
  extendedchars=true,              % lets you use non-ASCII characters; for 
%8-bits encodings only, does not work with UTF-8
  frame=single,                    % adds a frame around the code
  keepspaces=true,                 % keeps spaces in text, useful for keeping 
%indentation of code (possibly needs columns=flexible)
  keywordstyle=\color{blue},       % keyword style
  language=Octave,                 % the language of the code
  morekeywords={*,...},            % if you want to add more keywords to the set
  numbers=left,                    % where to put the line-numbers; possible 
%values are (none, left, right)
  numbersep=5pt,                   % how far the line-numbers are from the code
  numberstyle=\tiny\color{mygray}, % the style that is used for the line-numbers
  rulecolor=\color{black},         % if not set, the frame-color may be changed 
%on line-breaks within not-black text (e.g. comments (green here))
  showspaces=false,                % show spaces everywhere adding particular 
%underscores; it overrides 'showstringspaces'
  showstringspaces=false,          % underline spaces within strings only
  showtabs=false,                  % show tabs within strings adding particular 
%underscores
  stepnumber=2,                    % the step between two line-numbers. If it's 
%1, each line will be numbered
  stringstyle=\color{mymauve},     % string literal style
  tabsize=2,                       % sets default tabsize to 2 spaces
  title=\lstname                   % show the filename of files included with 
%\lstinputlisting; also try caption instead of title
}

%\title{EPSRC Centre for Doctoral Training in Medical Imaging\\\hfill\\MRes Project Plan}
% \title{Differential Diagnosis of Alzheimer subtypes through disease progression modelling}
% \author{
% Razvan Valentin Marinescu\\
% Student Number: 14060166\\
% \texttt{razvan.marinescu.14@ucl.ac.uk}
% }

\begin{document}
\belowdisplayskip=12pt plus 3pt minus 9pt
\belowdisplayshortskip=7pt plus 3pt minus 4pt


\begin{titlepage}
\begin{center}

% Upper part of the page. The '~' is needed because \\
% only works if a paragraph has started.
\includegraphics[width=0.2\textwidth]{ucl-logo2}~\\[1cm]

\textsc{\LARGE University College London}\\[1.5cm]

\textsc{\Large MRes Project Plan}\\[0.5cm]

\newcommand{\HRule}{\rule{\linewidth}{0.5mm}}

% Title
\HRule \\[0.4cm]
{ \Large Differential Diagnosis of Alzheimer subtypes using disease progression modelling \\[0.4cm] }

\HRule \\[1.5cm]

% Author and supervisor
\begin{minipage}{0.4\textwidth}
\begin{flushleft} \large
\emph{Author:}\\
R\u{a}zvan Valentin \textsc{Marinescu}\\
\href{razvan.marinescu.14@ucl.ac.uk}{razvan.marinescu.14@ucl.ac.uk}\\
\end{flushleft}
\end{minipage}
\begin{minipage}{0.4\textwidth}
\begin{flushright} \large
\emph{Supervisors:} \\
Prof. Daniel \textsc{Alexander}\\
\href{d.alexander@ucl.ac.uk}{d.alexander@ucl.ac.uk}\\[0.5cm]
Dr. Sebastian \textsc{Crutch}
\href{s.crutch@ucl.ac.uk}{s.crutch@ucl.ac.uk}
\end{flushright}
\end{minipage}

\vfill

EPSRC Centre for Doctoral Training in Medical Imaging\\ University College London

\vfill

% Bottom of the page
{\large \today}

\end{center}
\end{titlepage}
% \maketitle{}


\section*{Aims of the project}
% keep very short, just implementing an EBM model of PCA

We are planning on modelling PCA using an \textit{event-based model} (EBM), which models the progression of the disease as a sequence of events, where each event corresponds to a biomarker becoming abnormal. The biomarker dataset is multimodal and includes measurements of atrophy rates of cortical volumes, cognitive test scores and presence of proteins such as Amyloid beta.

%The EBM can calculate the most common progression of the events in the patient cohort and can be used to classify individual patients into the three main groups: cognitively normal, mild cognitive impairment and Alzheimer's disease. The EBM can also be used for longitudinal studies; for instance, it can stage patients after a predefined follow-up period.

\section*{Project Background}
% clinical challenges, PCA, why is the modelling challenging, other types of models.

Dementia and other related neurodegenerative disorders cost the UK approximately \pounds 26 billion a year. Developing a clear understanding of how these disorders operate and progress over time could help us find more effective treatment. This project is investigating the progression of Posterior Cortical Atrophy (PCA), which is an early onset variant of Alzheimer's disease that causes atrophy of the posterior part of the cerebral cortex, resulting in disruption of the visual and motor systems. 

Symptoms include blurred vision, impaired ability to read, difficulty with depth perception and navigating through space. The cause of PCA is still unknown and there are no fully accepted diagnostic criteria. PCA patients are often misdiagnosed with anxiety disorder or depression. Moreover, no specific and accepted scientific treatment for PCA has been found so far due to the rarity and variation of the disease.  All these factors together show that more research into PCA is needed and the EBM can assist to this. It could lead to better understanding of PCA progression, and can be used as a clinical tool for prognoses and differential diagnosis. 

\section*{Motivation}

The EBM has already been used to study familiar and sporadic Alzheimer's disease and Huntington's disease (HM Fonteijn et al, 2012, Young et al 2014).This project aims to apply these techniques in the study of PCA. We will be interested to find the most common event progression, use it to classify patients into several groups according to their disease stage. We will also do some longitudinal studies, such as finding the probability of a patient converting from cognitively normal to mild cognitive impairment or from mild cognitive impairment to PCA over time. The EBM would be very useful as a clinical tool because it can provide prognostic information regarding the progression of PCA. It could also be used to perform differential diagnosis, which is the aim of the project.


\section*{List of tasks \hfill Deadline}

\renewcommand\arraystretch{4.4} %\setlength\minrowclearance{2.4pt}
\newcommand\taskHeader[1]{\Large{\textbf{#1}}}

\subsection*{\noindent Data collection and preprocessing \hfill 10 April}

One needs to run segmentation algorithms to compute atrophy rates from MRI images of PCA patients. Data also needs to be collected and processed for other types of biomarker measurements, such as cognitive test scores and blood tests.

\subsection*{\noindent Implement the EBM \hfill 24 April}

One needs to read about the EBM and design the PCA-based EBM. The choice of distributions modelling the subject measurements might be different from previous EBM implementations.

\subsection*{\noindent Initial results: \hfill  8 May} 
% Positional variance matrices and EBM stage histograms  for different groups  
Use the EBM implementation to get some initial results for PCA: positional variance matrices and stage histograms for different groups.

\subsection*{\noindent Analyse follow-up data \hfill  15 May}
% , compute prob of remaining CN, MCI over time  
Analyse follow-up data and compute plots showing longitudinal consistency of patient disease stage. Also fit Cox proportional hazards models to find the probability of patients of converting from cognitively normal (CN) to mild cognitive impairment (MCI) or from MCI to PCA over time.

\subsection*{\noindent Build disease stage classifier \hfill  29 May}
% cut points, sensitivity and specificity  
Build a simple classifier using cut points in the histogram data which would be able to classify patients in the three main categories: CN, MCI and PCA.

\subsection*{\noindent Differential Diagnosis of PCA vs other AD subtypes \hfill 19 June}
Use the EBM to perform differential diagnosis of PCA and other Alzheimer's subtypes. 

\subsection*{\noindent Validation of the results \hfill  3 July}
Compare the results of the PCA-based EBM with other disease progression methods, such as those based on differential equations or self-modelling regression approaches.

\subsection*{\noindent Simultaneous MCMC sampling \hfill 20 July}

Implement MCMC sampling of both the model parameters and the most probable disease progression. This feature has not been tried in previous EBM implementations.

\subsection*{\noindent Write the MRes project report  \hfill  14 August}

Write the MRes project report showing all the results obtained so far.


% Temporal diagram hilighting the order in which areas of brain get affected by PCA
% (maybe, if relevant) Comparison of event sequences across left-right hemispheres?  \hfill  


\section*{Summary of progress to date}

I have so far experimented with the EBM by trying to reproduce some of the results of A. Young, 2012. We managed to get the MRI images from the Dementia Research Center. We are still missing the data from protein biomarkers and cognitive test scores. 


\end{document}
